\subsection{Algorithm pairs and what is held fixed}

Each pair consists of a \emph{faster-in-theory} algorithm $F$ and a
\emph{simpler} algorithm $S$ (Table~\ref{tab:pairs}). Three rules keep the
comparison about the algorithmic idea rather than about implementation effort.
\begin{enumerate}
  \item \textbf{Same language, compiler and flags.} C++17, GCC, \texttt{-O2},
  no \texttt{-march=native} (the x86-64 baseline ISA), no hand-written SIMD, no
  threads. (The harness can rebuild with \texttt{-O3 -march=native} for a
  secondary experiment; that experiment is not part of this paper.)
  \item \textbf{Same base case.} The recursive algorithms (Strassen, Karatsuba)
  bottom out in exactly the kernel they are compared against, so the measured
  difference is attributable to the recursive structure, its extra additions
  and its memory traffic.
  \item \textbf{Competent, textbook implementations on both sides.} Radix sort
  computes all digit histograms in one pass and skips passes in which every key
  shares a digit; Karatsuba uses a caller-provided scratch buffer; the
  Fibonacci heap pre-allocates one node per vertex. The simpler algorithms are
  what a careful programmer writes: \texttt{std::sort}, the $i$-$k$-$j$ loop
  order, 64-bit multiply-accumulate, an indexed binary heap with
  decrease-key---and, as a second baseline, the lazy \texttt{std::priority\_queue}
  idiom that is far more common in practice.
\end{enumerate}

\begin{table}[t]
\centering\small
\caption{The four algorithm pairs and the variables swept.}
\label{tab:pairs}
\begin{tabular}{p{2.2cm}p{3.3cm}p{3.3cm}p{5.6cm}}
\toprule
Pair & $F$ & $S$ & Swept \\
\midrule
Sorting & LSD radix sort, $\bigO(n\cdot 64/d)$ for $d$-bit digits & introsort (\texttt{std::sort}), $\bigO(n\log n)$ & $n\in[16,2^{26}]$; key width 8/16/32/64 bits; nearly-sorted, reverse, few-unique, organ-pipe; $d\in\{8,11,16\}$ \\
Matrix multiply & Strassen, $\bigO(n^{2.807})$ & naive triple loop, $\bigO(n^3)$ & $n\in[16,4096]$ incl.\ non-powers of two; cutoff $\in\{8,\dots,256\}$; $i$-$k$-$j$ vs $i$-$j$-$k$ baseline \\
Integer multiply & Karatsuba, $\bigO(n^{1.585})$ & schoolbook, $\bigO(n^2)$ & $n\in[8,2^{20}]$ 32-bit limbs; cutoff $\in\{8,\dots,128\}$; random and all-ones limbs \\
Shortest paths & Dijkstra + Fibonacci heap, $\bigO(m+n\log n)$ & Dijkstra + binary heap, $\bigO((n+m)\log n)$ & $n\in[10^2,8.4\times10^6]$; average degree $\in\{2,\dots,256\}$; random / 2-D grid / complete graphs; lazy baseline \\
\bottomrule
\end{tabular}
\end{table}

\subsection{Inputs}
Inputs are generated deterministically from a 64-bit seed combined with $n$
(xoshiro256** seeded through splitmix64), so every cell of the experiment is
reproducible from the command line that produced it. Sizes form a geometric
grid (factor 2 for the main sweeps, 1.5 for integer multiplication, and a
hand-picked list for matrices that includes sizes such as 24, 48 and 96 that
exercise Strassen's padding path). A \emph{cell} is a (pair, input family,
parameter, $n$) tuple. Cells that make no sense are skipped by rule (a complete
graph on $10^7$ vertices); an algorithm whose single run exceeds a wall-clock
budget is dropped for larger $n$ in that series, so that cubic baselines do not
dominate the running time of the experiment.

\subsection{Measurement}
\paragraph{Timer.} \texttt{std::chrono::steady\_clock} (\texttt{CLOCK\_MONOTONIC}
through the vDSO, about 20\,ns overhead). For inputs whose single run exceeds
2\,ms the input is copied outside the timed region and only the algorithm is
timed. Shorter runs are batched: the harness runs $K$ iterations of (copy
input; run algorithm) and reports the per-iteration mean, with $K$ chosen from
a warm calibration run; because in-place algorithms need a fresh copy per
iteration, the copy is inside the timed region, and its cost is measured
separately with the same $K$ on the same pre-allocated buffer and subtracted
in the analysis (never below 5\% of the raw time).

\paragraph{Repetitions and interleaving.} Each cell is measured $r=5$ times
after $w=2$ untimed warm-up runs. Within the repetition loop the algorithms
alternate in a rotating order, so slow drift in machine state (thermal
throttling, background activity, frequency changes) is shared equally by $F$
and $S$. The per-cell statistic is the median; the interquartile range is
retained and plotted.

\paragraph{Hardware counters.} When \texttt{perf\_event\_open} is permitted,
each measurement also records cycles, retired instructions, L1d read misses,
last-level-cache misses and branch mispredictions for the user-space portion of
the run, scaled for multiplexing. These are explanatory, not primary: the
crossover is defined on wall-clock time. Where counters are unavailable (virtual
machines, WSL2) the study is timing-only and says so.

\paragraph{Correctness.} Every benchmark binary cross-checks the outputs of $F$
and $S$ on the last repetition of each cell and aborts on mismatch. Separate
randomized test suites compare each implementation with an independent
reference (\texttt{std::sort}; the $i$-$j$-$k$ product; schoolbook plus the
modular identity $(ab)\bmod p = ((a\bmod p)(b\bmod p))\bmod p$ with
$p=2^{61}-1$; Bellman--Ford), and both heaps are exercised directly against a
scan-for-minimum oracle under random push, decrease-key and pop sequences.

\subsection{Runs and a measurement correction}\label{sec:method-patch}
The reported dataset is the union of three runs made in one session on the
same machine, with the same binary flags: \texttt{full} (the sweeps of
Table~\ref{tab:pairs}, 5.3 hours), \texttt{refine} (factor-1.19 grids over
the neighbourhood of each crossing, 20 minutes) and \texttt{patch} (70
minutes). The patch exists because the first version of the harness measured
the input-copy baseline into a freshly allocated buffer, so that for large
inputs the baseline included the operating system's page-fault cost while the
copy inside the timed loop, reusing the buffer, did not; the subtraction then
removed too much. The effect is below 5\% wherever the copy is small relative
to the run, which is almost everywhere, but for complete graphs (470\,MB of
adjacency at 7\,700 vertices against a 100\,ms run) and for the largest
arrays sorted by the fastest configurations it was large enough to invalidate
the cell. We fixed the harness as described above, re-measured every series in
which the copy had exceeded 5\% of the raw time (dense graphs at all sizes,
arrays of $2^{21}$ keys and above, sparse graphs of $2^{20}$ vertices and
above, and a handful of small cells), and let the re-measurement
\emph{replace} the original cells over the range of $n$ it covers rather than
pool with them. Both the original and the corrected data are in the
repository; the analysis script's \texttt{--supersede} flag applies the
replacement.

\subsection{Machine}
All reported measurements were made on a laptop with an Intel Core
i7-14650HX (Raptor Lake, 8 performance cores + 8 efficiency cores, 24
threads; 48\,KB L1d and 2\,MB L2 per performance core, 30\,MB shared L3)
running Ubuntu under WSL2 (Linux 6.18, Microsoft hypervisor) with 7.6\,GB of
memory visible to the guest, GCC 15.2.0 with \texttt{-O2}, CMake 4.2. The
laptop was on mains power in the manufacturer's \emph{Performance} mode, the
Windows power plan was set to \emph{High performance} with maximum processor
state 99\% (which disables Turbo Boost on Intel parts), the WSL virtual
machine was restricted to the performance cores through its process
affinity, and the benchmark thread was pinned to one virtual CPU. The hypervisor hides
hardware performance counters, so the reported run is timing-only; the
harness collects counters automatically on bare-metal Linux, where it also
sets the performance governor, disables turbo and pins the frequency
(\texttt{scripts/setup\_machine.sh}). We report the machine exactly and treat
absolute times as specific to it; the crossover \emph{ordering} and
\emph{order of magnitude} are the transferable results.

\subsection{Crossover estimation}
For a series (pair, family, parameter) let $\rho(n)=\tilde t_F(n)/\tilde t_S(n)$
be the ratio of median runtimes on the size grid. The crossover $n^*$ is the
\emph{sustained} crossing: the grid interval after which $\rho<1$ at every
larger measured size, interpolated linearly in $(\log n,\log\rho)$. A
transient dip below 1 followed by a return above it is not counted; it remains
visible in the ratio plots and in the reported minimum ratio. Series with $\rho<1$ already at the smallest $n$ are
reported as \emph{always faster} (from $n_{\min}$) and series with $\rho\ge1$
throughout as \emph{never faster} (up to $n_{\max}$), together with
$\rho(n_{\max})$ so the reader can see how far from crossing they are.
Uncertainty comes from a non-parametric bootstrap: the $r$ repetitions of every
cell are resampled with replacement (1000 replicates), medians and $n^*$ are
recomputed, and the shortest interval containing 95\% of the replicates (in
$\log n$) is reported. The shortest rather than the central interval is used
because the sustained-crossing statistic can be bimodal: a resample in which
one noisy mid-range cell pops above $\rho=1$ relocates the crossing to the far
end of the sweep, and a percentile interval would stretch over that tail even
when nearly all of the mass sits tightly around the true crossing.

Two caveats govern how the intervals should be read. First, $n^*$ is
interpolated between grid points, so an interval cannot resolve a crossing
more finely than the grid allows; the main sweeps use a factor-2 grid and the
\texttt{refine} profile re-sweeps the neighbourhood of each crossing on a
factor-1.19 grid for that purpose. Second, and more important, the width of an
interval is itself a result. We classify each crossing as \emph{sharp} when
its interval is narrower than one grid step or than a factor of 1.5,
whichever is wider (a $\pm20\%$ interval is sharp by any practical standard,
even on a fine grid), as \emph{gradual} otherwise, and as \emph{marginal}
when the faster algorithm is never more than 10\% ahead by the largest size
measured, whatever the interval. A gradual crossing means that $\rho$ approaches
1 slowly, i.e.\ the two algorithms are within measurement noise of each other
across the whole interval; the honest statement is then not ``the fast
algorithm wins from $n^*$'' but ``over this range it makes no measurable
difference, and beyond it the fast algorithm wins by the ratio we report at
$n_{\max}$''. Gradual crossings are reported as ranges, never as points. Empirical scaling exponents are fitted as $\log t = b\log n + c$
over the largest decade of $n$ for each algorithm, to show how far the
measured growth is from the asymptotic bound in the measured range.
